% ============================================================================
% ICSE Paper — Approach Section
% ILinker2: LLM-Driven Trace Link Recovery between Architecture Documents and Models
% ============================================================================

\section{Approach}\label{sec:approach}

\subsection{Overview}\label{sec:overview}

We propose an LLM-driven pipeline for recovering trace links between
Software Architecture Documents (SAD) and Software Architecture Models (SAM).
Given a natural-language architecture document~$D$ consisting of sentences
$\{s_1, \ldots, s_n\}$ and a component model~$M$ containing components
$\{c_1, \ldots, c_m\}$, our approach produces a set of trace links
$L \subseteq S \times C$, where each link $(s_i, c_j)$ asserts that
sentence~$s_i$ discusses component~$c_j$.

The pipeline is organized as a \emph{fan-out / fan-in} staged DAG
(Figure~\ref{fig:pipeline}) with three tiers:

\begin{enumerate}
    \item \textbf{Knowledge acquisition} (Tier~1): Independent analyses
      extract structured knowledge from the raw inputs---model structure,
      document vocabulary, and linguistic patterns.
      Model analysis, document analysis, and seed extraction share
      no data dependencies; pattern learning and multiword enrichment
      depend on model analysis results and run in a brief enrichment
      step once the initial analyses complete.
    \item \textbf{Link recovery} (Tier~2): Four complementary strategies
      produce candidate links via different mechanisms---high-precision seed
      extraction, explicit entity extraction, pronominal coreference
      resolution, and deterministic partial-reference injection. The
      strategies are largely independent: each targets a different type
      of reference, and their outputs converge only at the merge point.
    \item \textbf{Merge and filtering} (Tier~3): A single fan-in stage
      merges all candidate links, applies convention-aware boundary filtering
      to \emph{all} links regardless of provenance, and removes residual
      false positives via a lightweight deterministic evidence filter.
\end{enumerate}

\noindent This architecture offers two key advantages.
First, the recovery strategies contribute additively---seed extraction
captures high-precision explicit links, entity extraction finds remaining
mentions, coreference recovers pronoun-based references, and partial
injection catches abbreviated names---without interference or ordering
constraints between strategies.
Second, the separation of knowledge acquisition from link recovery
enables each knowledge source to be refined or replaced independently
without affecting the recovery strategies.

% ---------------------------------------------------------------------------
\begin{figure}[t]
\centering
\small
\begin{verbatim}
    Tier 1: Knowledge Acquisition (concurrent)
  ┌──────────┐  ┌──────────┐  ┌──────────┐
  │  Model    │  │ Document │  │   Seed   │
  │ Analysis  │  │ Analysis │  │Extraction│
  └─────┬────┘  └────┬─────┘  └────┬─────┘
        │             │             │
        ├─────────────┤             │
        ▼             │             │
  ┌──────────┐        │             │
  │ Pattern  │        │             │
  │ Learning │        │             │
  └─────┬────┘        │             │
        └──────┬──────┘             │
               │  shared knowledge  │
               ▼                    │
  ┌──────────────┐ ┌───────────┐   │
  │   Entity     │ │Coreference│   │
  │  Extraction  │ │Resolution │   │
  │+ Validation  │ │           │   │
  └──────┬───────┘ └─────┬─────┘   │
         │               │         │
         ├───────────────┤         │
         ▼               │         │
  ┌──────────┐           │         │
  │ Partial  │           │         │
  │Injection │           │         │
  └────┬─────┘           │         │
       │   Tier 2: Link Recovery   │
       └──────────┬───────────────┘
                  ▼
        ┌──────────────────┐
        │ Merge + Boundary │
        │ Filter + Evidence│
        │     Filter       │
        └──────────────────┘
       Tier 3: Merge & Filter
\end{verbatim}
\caption{Fan-out / fan-in pipeline architecture. Tier~1 acquires knowledge
concurrently; pattern learning depends on model analysis and runs after the
initial analyses. Tier~2 runs link recovery strategies that converge at the
merge point. Seed extraction requires no knowledge inputs and runs in Tier~1.}
\label{fig:pipeline}
\end{figure}
% ---------------------------------------------------------------------------

\subsection{Knowledge Acquisition}\label{sec:knowledge}

Tier~1 builds a shared knowledge base from the raw document and model.
Model analysis, document analysis, and seed extraction are mutually
independent---each reads only the raw inputs---enabling concurrent
execution. Pattern learning and multiword enrichment depend on the
model analysis results (filtered component list and generic-word
classification) and run once the initial analyses complete.

\subsubsection{Model Analysis}\label{sec:model-analysis}

We analyze component names from the architecture model to distinguish
\emph{architectural} names from \emph{ambiguous} ones.
Architectural names (e.g., \texttt{TaskScheduler}, \texttt{RenderEngine})
identify specific system roles, while ambiguous names
(e.g., \texttt{Core}, \texttt{Adapter}) are single common-English words
that documentation authors frequently use generically.

Classification uses a few-shot LLM prompt with examples drawn from
software-engineering textbook domains (compilers, operating systems,
networking), followed by a deterministic structural guard:
multi-word names, CamelCase compounds (detected via the regex pattern
\texttt{[a-z][A-Z]}), and all-uppercase acronyms are always classified
as architectural, regardless of LLM output.
This hybrid approach leverages LLM judgment for ambiguous cases while
guaranteeing that structurally unambiguous identifiers are never
misclassified.

The analysis additionally extracts: (a)~\emph{implementation relationships}
between component names (e.g., \texttt{BasicScheduler}~$\to$~\texttt{Scheduler}),
identified by shared substrings; (b)~\emph{shared vocabulary} across
components, mapping words that appear in multiple component names;
and (c)~\emph{generic partial words}, the set of ambiguous-classified
words that also appear as subwords of other component names, used later
to calibrate partial-reference matching.

\subsubsection{Document Analysis}\label{sec:doc-analysis}

Document analysis extracts two kinds of information: structural
characteristics and alternative component names.

\paragraph{Document profiling.}
We compute the pronoun ratio (fraction of sentences containing anaphoric
pronouns such as \emph{it}, \emph{they}, \emph{this}), component mention
density, and sentence-to-component ratio.
These metrics are logged for diagnostic purposes and inform threshold
calibration in downstream phases.

\paragraph{Alternative name extraction.}
An LLM extracts three categories of alternative names:
\begin{itemize}
    \item \emph{Abbreviations}: short forms introduced via parenthetical
      definitions (e.g., ``Abstract Syntax Tree (AST)'' defines AST).
    \item \emph{Synonyms}: alternative names that unambiguously identify
      exactly one component.
    \item \emph{Partial references}: trailing words of multi-word component
      names used alone (e.g., ``Dispatcher'' for \texttt{EventDispatcher}).
\end{itemize}

\noindent A calibrated few-shot judge validates each proposed mapping.
The judge auto-approves abbreviations formed from component-name initials,
trailing words of multi-word names, and CamelCase identifiers, while
rejecting generic terms (\emph{system}, \emph{process}).
A deterministic \emph{CamelCase rescue} overrides any judge rejection of
CamelCase identifiers, since constructed identifiers are never generic
English words.

An enrichment step further expands partial references: for each multi-word
component, if its trailing word appears in at least three sentences where
the full name is absent, and no other component shares that trailing word,
the partial mapping is added automatically.
For trailing words that are generic (i.e., appear in the ambiguous set),
the match requires capitalized occurrences to avoid false positives.

A subsequent \emph{partial usage classification} step uses per-partial
LLM calls to determine whether each single-word generic partial is
typically used as a component reference (NAME) or as an ordinary English
word (ORDINARY) in the document.
ORDINARY-classified partials are excluded from the confirmed-alias
context provided to the Tier~3 boundary filter, preventing the filter
from treating generic uses as validated component references.

\subsubsection{Pattern Learning}\label{sec:patterns}

Two independent LLM passes identify \emph{subprocess terms}---words that
refer to internal parts of components rather than the components themselves
(e.g., ``garbage collector'' inside a \texttt{MemoryManager}).
The first pass proposes candidates; the second cross-validates them in a
debate structure, rejecting terms that may be valid component references.
This debate mechanism reduces the risk of over-filtering in later
validation phases.
The phase also extracts \emph{action indicators} (verbs used when a
component performs an action) and \emph{effect indicators} (verbs
describing results), which are used as contextual signals during
validation.

Pattern learning reads the filtered component list (excluding
implementation names identified by model analysis) and the document
sentences.
This dependency places pattern learning after model analysis completes,
forming a brief enrichment step between knowledge acquisition and link
recovery.

% ---------------------------------------------------------------------------

\subsection{Link Recovery Strategies}\label{sec:recovery}

Tier~2 comprises four complementary strategies that produce candidate links.
Each targets a different \emph{type} of reference---high-precision seed
links, explicit entity mentions, pronominal references, and abbreviated
names---and their contributions converge only at the Tier~3 merge point.
Seed extraction requires no knowledge inputs and runs concurrently
with Tier~1 knowledge acquisition; the remaining three strategies
consume the shared knowledge base.

\subsubsection{Seed Extraction}\label{sec:seed}

The pipeline accepts either an external baseline (e.g., TransArc CSV)
or runs a self-contained LLM-based seed extractor called \textbf{ILinker2}.
ILinker2 is a high-precision explicit link extractor that
operates independently of all knowledge phases---its only inputs are the
raw sentences and component list.

ILinker2 uses two complementary LLM passes over batched document segments
(50~sentences per batch with 5-sentence overlap):

\begin{itemize}
    \item \textbf{Pass~A (extraction-framed):} ``For each sentence, identify
      which components are explicitly mentioned.'' This pass is optimized
      for exhaustive recall of named references.
    \item \textbf{Pass~B (actor-framed):} ``For each sentence, determine
      which component is the subject or primary actor.'' This reframing
      surfaces references that Pass~A may miss by focusing on agency
      rather than name matching.
\end{itemize}

\noindent Both passes share exclusion rules that reject names appearing
only inside dotted package paths (e.g., \texttt{renderer.utils.config}
does not reference \texttt{Renderer}) and generic English words used in
their ordinary sense.

The merge strategy applies \emph{asymmetric voting}:
exact-name matches from \emph{either} pass are accepted (union), while
non-exact matches (synonyms, partials) require agreement from \emph{both}
passes (intersection).
This exploits the observation that exact matches are reliably identified
by a single pass, while fuzzy matches benefit from independent agreement
to suppress hallucinations.

ILinker2 achieves approximately 95\% precision but only 76\% recall as a
standalone extractor, deliberately trading recall for precision since
the remaining strategies recover the missing links.

\subsubsection{Entity Extraction and Validation}\label{sec:entity}

\paragraph{Extraction.}
An LLM prompt, enriched with the abbreviations, synonyms, and partial
references from document analysis, identifies all component references in
batches of 50~sentences.
The prompt covers six inclusion patterns: direct name matches,
space-separated forms of compound names, functional descriptions,
known aliases, interaction participants (``X sends data to Y'' references
both X and Y), and passive/prepositional mentions (``handled by X'').
An \emph{abbreviation guard} rejects invalid abbreviation expansions by
verifying that the context following the abbreviation is consistent with
the full component name.

For components that receive no links from either the seed or batch
extraction, a \emph{targeted recovery} step issues single-component
prompts with all known aliases, recovering references missed by the
broad-scope batch extraction.

\paragraph{Validation.}
Candidates undergo a three-step validation process:

\begin{enumerate}
    \item \emph{Code-first auto-approval:} Candidates where a component
      name or known alias appears verbatim in the sentence text (verified
      by word-boundary matching in code) are approved without LLM cost.
    \item \emph{Two-pass LLM intersection:} Remaining candidates are
      validated by two independent LLM passes with complementary
      foci---one on \emph{actor role} (``is the component performing
      an action?'') and one on \emph{direct reference} (``does the text
      refer to the specific architectural component?'').
      A candidate is accepted only if \emph{both} passes approve it,
      applying intersection voting for precision.
    \item \emph{Evidence post-filter:} For components classified as
      ambiguous by model analysis, an additional LLM pass must cite the
      \emph{exact substring} from the sentence that names the component.
      The cited text is verified against the actual sentence content in
      code, rejecting hallucinated evidence.
\end{enumerate}

\subsubsection{Coreference Resolution}\label{sec:coref}

Sentences containing anaphoric expressions---pronouns (\emph{it},
\emph{its}, \emph{they}, \emph{their}, \emph{this}, \emph{these},
\emph{that}, \emph{those}) and definite descriptions (\emph{the
component}, \emph{the service})---are analyzed to resolve references
back to architecture components.

The pipeline uses a unified \emph{cases-in-context} strategy:
each pronoun-bearing sentence is presented as an individual case with
a $\pm5$-sentence bidirectional context window, processed in batches
of~10.
The LLM sees each case's surrounding context (marked with
\texttt{>>>} for the target sentence) and resolves any pronoun
that refers to an architecture component.
This per-case presentation with local context is robust across both
simple and complex documents, eliminating the need for document
complexity gating.

The resolution enforces a \emph{citation-grounded verification}
protocol.
The LLM must cite an antecedent sentence where the component was
explicitly named, along with the exact quote containing the component
name.
This citation is then verified deterministically in code:
\begin{enumerate}
    \item The component name (or a known alias) must appear verbatim in
      the cited antecedent sentence.
    \item The antecedent must be within 3~sentences of the pronoun.
    \item Sentences identified as subprocess descriptions by pattern
      learning are excluded.
\end{enumerate}

\noindent This hybrid LLM-plus-code verification eliminates a class of
hallucinated coreference links where the LLM fabricates a plausible but
nonexistent antecedent.
Because the antecedent verification is strict (verbatim name match
within 3~sentences), false positives from hallucinated resolutions are
effectively eliminated at the code level, making coreference links
safe to accept without further LLM review in Tier~3.

\subsubsection{Partial Reference Injection}\label{sec:partial}

For each partial reference learned during document analysis
(e.g., ``Dispatcher''~$\to$~\texttt{EventDispatcher}), the pipeline
deterministically scans all sentences for word-boundary matches, injecting
links for any sentence-component pairs not already covered by the seed,
entity extraction, or coreference strategies.
A \emph{clean mention} filter ensures the partial name appears as a
standalone word---not inside a dotted path (\texttt{x.Dispatcher.init}),
a hyphenated compound, or adjacent to a period indicating a qualified name.

This strategy is entirely deterministic (no LLM calls), providing
stable recall for abbreviated references at zero additional inference cost.
Its output is filtered by the convention-aware boundary filter in Tier~3.

% ---------------------------------------------------------------------------

\subsection{Merge and Filtering}\label{sec:merge}

Tier~3 merges the outputs of all recovery strategies and applies
two layers of filtering to remove false positives.

\subsubsection{Priority-Based Merge}\label{sec:dedup}

Links from all sources (seed, entity extraction, coreference, partial
injection) are deduplicated by $(sentence, component)$ pair.
When multiple sources produce the same link, the highest-priority source
is retained according to a fixed ranking:
$\text{seed} > \text{validated} > \text{entity} > \text{coreference} > \text{partial\_inject}$.
An optional \emph{parent-overlap guard} removes links to implementation
components when the abstract parent component is already linked to the
same sentence (e.g., if both \texttt{Scheduler} and
\texttt{BasicScheduler} are linked to a sentence, the implementation
variant is dropped).

\subsubsection{Convention-Aware Boundary Filter}\label{sec:boundary}

A structured LLM filter applies a 3-step reasoning guide to \emph{all}
links, regardless of provenance (including seed links):

\begin{enumerate}
    \item \textbf{Hierarchical name reference (Step~1):}
      Reject if the component name appears only as a prefix in a
      dotted or qualified path (e.g., ``\texttt{X.utils} provides
      helper functions'' does not reference component~X).
      An exception preserves links when the sentence also names
      the component with the word ``component''
      (e.g., ``for the X component'').
    \item \textbf{Entity confusion (Step~2):}
      Reject if the name is used as a technology, methodology, or
      generic English word in a non-architectural sense.
      A critical exception preserves links when a component is
      \emph{named after} a technology (e.g., a component called
      ``Kafka Broker'')---any sentence describing that technology's
      capabilities is about the component.
    \item \textbf{Default (Step~3):}
      If neither rejection pattern applies, the link is kept.
\end{enumerate}

\noindent Structural guardrails ensure multi-word component names and
CamelCase identifiers are never treated as generic words, and that
interaction sentences produce links for all participating components.

\subsubsection{Evidence Filter}\label{sec:evidence-filter}

After the boundary filter, a lightweight deterministic
\emph{evidence filter} removes links that lack textual grounding
in the sentence.
A link $s \to c$ is kept if \emph{any} of the following holds:

\begin{enumerate}
    \item The component name appears as a standalone word-boundary
      match in the sentence text (case-sensitive for single-word names,
      case-insensitive for multi-word names), excluding occurrences
      inside dotted paths or hyphenated compounds.
    \item A known alias (synonym, abbreviation, or partial reference
      from document analysis) appears in the sentence text.
    \item The link originates from the seed extractor, which already
      verified explicit mentions.
    \item The link originates from coreference resolution, which
      already verified the antecedent citation.
\end{enumerate}

\noindent Links that satisfy none of these conditions---where the
component name is absent from the sentence and no alias or prior
verification provides evidence---are dropped.
This filter requires zero LLM calls and removes residual false
positives introduced by entity extraction or validation when the
LLM hallucinated a reference that does not appear in the text.

The rationale for trusting coreference links without additional
review is that antecedent verification (\S\ref{sec:coref}) already
provides strong evidence: the component was explicitly named in a
nearby sentence and the pronoun grammatically refers back to it.
Empirically, applying an LLM judge to coreference links kills true
positives but never catches false positives, since the code-level
verification is more reliable than LLM re-evaluation.

% ---------------------------------------------------------------------------

\subsection{Cross-Cutting Design Principles}\label{sec:principles}

Three principles guide the pipeline design across all tiers:

\paragraph{Asymmetric voting.}
We use \emph{intersection} voting (both passes must agree) for
precision-critical decisions: seed merge for non-exact matches
(\S\ref{sec:seed}) and entity validation (\S\ref{sec:entity}).
This reflects the observation that false positives in early stages
propagate through the pipeline and are difficult to remove later.

\paragraph{Citation-grounded verification.}
Rather than trusting LLM judgments as opaque decisions, we require the
LLM to produce \emph{verifiable evidence}---antecedent sentence citations
in coreference, exact-substring evidence in validation---that is checked
deterministically in code.
This hybrid approach leverages LLM reasoning for semantic understanding
while using programmatic guards to catch hallucinations.
The pattern applies in two places: coreference antecedent verification
(\S\ref{sec:coref}) and the evidence post-filter for ambiguous names
(\S\ref{sec:entity}).
A complementary but purely deterministic guard---the abbreviation guard
(\S\ref{sec:entity})---validates abbreviation expansions via string
matching without LLM involvement.

\paragraph{Benchmark-clean prompt design.}
All few-shot examples, reasoning guides, and prompt templates use
abstract placeholders (X, Y) or safe software-engineering textbook
domains (compilers, operating systems, networking).
No benchmark component names, dataset-specific vocabulary, or
gold-standard information appears in any prompt, ensuring that
measured performance reflects genuine generalization rather than
data leakage.
